\section{Introduction.}
%
\begin{print}
 The SCFE project was developed for solution of partial differential equations by means of finite element method. It was projected, keeping in mind to use it for the implementation of adaptivity for finite element methods. Another important criteria for design of this code is ability to solve coupled field problems. Code of SCFE was implemented on C++ programming languages due to object-oriented concept of it. Hierarchy of classes, virtual classes, mechanizmes of templates play an important role throughtout implementation. For the code development we used GNU C++ compiler( version 2.95)  on LINUX and CC on SGI. At present scalar problem such as acoustic equation is implemented, which is solved by using the Newmark method.\\
 I will give a sketch of design of programm and some issues of implementation. 
\end{print}
%
\newslide
\newpage
%
\section{ Design features .}
\begin{itemize}
%
\item{\em separation of topological and algebraic entities}\\
%
\begin{print}
 It is important criteria for design of code for implementation of adaptivity due to complex of work with grid. By using global node-numbering strategy we are able to transfer structure of grid into algebraic structure.
\end{print}
%
\item{\em clear structure of code}\\
\begin{print}
 It should be easy to catch structure of programm for users to be able to implement its own problems. So it should be desighned the way to hide and shield some structures and data from users.
\end{print}
%
\item{\em efficiency}\\
\begin{print}
 At the same time it should be efficient, because computations are often time-consuming, especially in 3D. Of course, this and former criteria are contracdictory and some time must be traded against each other.
\end{print}
%
\item{ \em features of coupled field problems }\\
\begin{print}
 ???????
\end{print}
\end{itemize}
%
\newslide
%
\section{ Structure of code.}
\begin{figure}[H]
\caption{ Structure of SCFE }
\begin{center}
%\includegraphics [width=8cm, angle=0] {structure}
\end{center}
\end{figure}
\newslide
\begin{print}
 The main classes of code are shown in figure. I will do a short description of
these classes.
\end{print}

